% Generated draft: requires mathematical and source review.
\documentclass[11pt]{article}
\usepackage[margin=25mm]{geometry}
\usepackage{fontspec}
\setmainfont{Latin Modern Roman}
\usepackage{amsmath,amssymb,mathtools,unicode-math}
\setmathfont{Latin Modern Math}
\usepackage{enumitem,needspace,xurl,hyperref}
\setlength{\parindent}{0pt}
\setlength{\parskip}{5pt}
\setlength{\emergencystretch}{3em}
\newcommand{\N}{\mathbb{N}}
\newcommand{\Z}{\mathbb{Z}}
\newcommand{\Q}{\mathbb{Q}}
\newcommand{\R}{\mathbb{R}}
\newcommand{\C}{\mathbb{C}}
\newcommand{\F}{\mathbb{F}}
\newcommand{\A}{\mathbb{A}}
\newcommand{\PP}{\mathbb{P}}
\newcommand{\E}{\mathbb{E}}
\newcommand{\eps}{\varepsilon}
\newcommand{\dd}{\,\mathrm{d}}
\begin{document}
\section*{UnsolvedMath ID: 6000002}
\textbf{Source problem code:} catalogue entry
\subsection{Definitions and mathematical statement}
For a smooth manifold $M$, write $\mathcal D_+(M)$ for the smooth strictly positive probability densities: smooth positive sections of the density bundle whose total integral is one. Relative to a fixed positive reference density $\mu_0$, these are $\rho\mu_0$ with
\[
\rho\in C^\infty(M),\qquad \rho>0,\qquad \int_M\rho\,\mathrm d\mu_0=1.
\]
For $\R^n$ the usual reference is Lebesgue measure and integrability must be imposed; smooth positivity alone is insufficient.

The supplied premise says that $\mathcal D_+(S^1)$ is diffeomorphic to the manifold of smooth equiaffine transformations of $S^1$. The problem asks for corresponding geometric descriptions of $\mathcal D_+(S^n)$, particularly for $n\geq2$, and $\mathcal D_+(\R^n)$.

The record does not define ``smooth equiaffine transformations'' of a circle or specify the higher-dimensional transformation spaces to be sought. In its ordinary ambient finite-dimensional meaning, an equiaffine transformation is an affine map $x\mapsto Ax+b$ with determinant-one linear part; such a finite-dimensional group is not the infinite-dimensional manifold of all smooth positive densities. Thus that ordinary meaning cannot simply be substituted for the source's unspecified circle construction. The action, normalization/stabilizer, and intended infinite-dimensional smooth structure must be recovered from the cited context. For $\R^n$, topology, decay, and behavior at infinity require additional choices. No definite diffeomorphism theorem for the proposed higher-dimensional spaces is stated without these data.
\subsection{Short English statement}
Find geometric transformation-space models for smooth positive probability densities on spheres and Euclidean spaces, extending the source's circle identification. The intended equiaffine transformation space and infinite-dimensional smooth structure are not defined.
\subsection{Sources}
\begin{itemize}
\item Local UnsolvedMath registry, permanent ID 6000002, exact source code catalogue entry, ``Probability Densities and Equiaffine Transformations''; proposed by Shun-ichi Amari, 1998. Registry status partially solved, classification PARTIAL-PROGRESS.
\item Furuhata, Matsuzoe, and Urakawa, \emph{Open Problems in Affine Differential Geometry and Related Topics} (1998), item 1(b), scanned PDF pages 1--2, \url{https://www.jstage.jst.go.jp/article/iis/4/2/4_2_125/_pdf/-char/en}. Supplied extraction: manual OCR from page images, accessed July 29, 2026. No PDF or page image was accessed for this draft.
\item The supplied research summary calls the circle case classical, mentions diffeomorphism-group/projective interpretations, and says the higher-sphere and Euclidean versions appear open or not cleanly established. It suggests possible reinterpretation using full diffeomorphism groups or density bundles but supplies no precise reference or definition. These are interpretive supplied assessments, not a verified identification of the intended spaces.
\item Provenance says the article presents the item as open and flags current status NEEDS\_REVIEW. The supplied rights note says the J-STAGE article page marks the article CC BY 4.0; attribution is retained here.
\end{itemize}
\end{document}
